\section*{Table 1. Characteristics of the active study pool}

\begin{longtable}{p{0.16\textwidth} p{0.25\textwidth} p{0.20\textwidth} p{0.29\textwidth}}
\toprule
Study & Occupational population & Design & Main exposure/intervention \\
\midrule
\endfirsthead
\toprule
Study & Occupational population & Design & Main exposure/intervention \\
\midrule
\endhead
    Alavi (2015) & Office workers & Analytical cross-sectional & Lifestyle/occupational factors \\
    Garbarino (2015) & Police officers & Prospective cohort & Work stress \\
    Kuwahara (2016) & Japanese workers & Prospective cohort & Physical activity \\
    Ko (2016) & Male public-office workers & Analytical cross-sectional & Physical activity \\
    Browne et al. (2017) & Sedentary occupation workers & Analytical cross-sectional & Physical activity \\
    Yamaguchi (2018) & Japanese manufacturing workers & Prospective cohort & Work-related stress \\
    Haufe (2019) & Volkswagen company employees & Randomized controlled trial & Exercise \\
    Appiah (2020) & Commercial taxi drivers & Analytical cross-sectional & Diet/lifestyle/physical inactivity \\
    Aboma (2020) & Working adults & Analytical cross-sectional & Lifestyle/diet/sedentary behaviour \\
    Jovanovic et al. (2020) & Security guards & Analytical cross-sectional & Occupational stress \\
    van Zon et al. (2020) & Dutch workers across occupational groups & Cohort/longitudinal observational & Occupation/physical activity/diet \\
    van Zon et al. (2021) & Older workers & Prospective cohort & Diet/physical activity/health behaviours \\
    Eftekhari (2021) & Medical university staff & Analytical cross-sectional & Job stress \\
    Apprey (2022) & Fuel tanker drivers & Analytical cross-sectional & Dietary intake \\
    Zhang (2024) & Petrochemical workers & Analytical cross-sectional & Occupational stress \\
    Sharifi (2025) & Health workers & Analytical cross-sectional & Diet quality \\
    Ma (2025) & Occupational men at oilfield canteens & Randomized controlled trial & Dietary intervention \\
    dos Reis et al. (2026) & Healthcare workers & Analytical cross-sectional & Night work/diet \\
    Petek \& Karabudak (2026) & Heavy-industry workers & Analytical cross-sectional & Diet/nutritional status \\
    Reinoso-Barbero et al. (2026) & Multinational-company employees & Analytical cross-sectional & Mediterranean diet \\
    Ayeltigah et al. (healthcare) (2026) & Healthcare professionals & Analytical cross-sectional & Stress/diet/physical activity \\
    Ayeltigah et al. (university) (2026) & University workers & Analytical cross-sectional & Stress/diet/physical activity \\
    Ayeltigah et al. (military) (2026) & Military personnel & Analytical cross-sectional & Work stress/diet/physical activity \\
    Earnest (2015) & Employees from 93 companies & Quasi-experimental pre-post & Worksite behavioral weight-loss intervention \\
    Kim et al. (2016) & Male white-collar employees & Quasi-experimental pre-post & Dietary intervention and worksite food environment \\
    Ryu (2017) & Office workers & Quasi-experimental pre-post & Workplace health promotion: education/self-monitoring/diet/activity \\
    So (2025) & Employees & Quasi-experimental single-group longitudinal & HIIT plus calorie restriction \\
    Bogale (2025) & Bank employees/office workers & Randomized controlled trial & Lifestyle education: diet/physical activity/stress and other risk factors \\
    Kim et al. BEST (2015) & Korean male workers with MetS & Controlled nonrandomized trial & Physical activity/diet/counselling \\
    Choi (2017) & Middle-aged male office workers & Randomized controlled trial & T'ai Chi + health education \\
    Huang et al. (2017) & Taiwanese workers & Analytical cross-sectional & Occupational/leisure/commuting physical activity \\
    Fang et al. (2019) & Overweight sedentary high-tech employees & Randomized controlled trial & Physical activity program \\
    Schilling et al. (2020) & Police officers & Prospective cohort & Work stress/physical activity/cardiorespiratory fitness \\
    Park \& Hwang (2024) & Male office workers with MetS & Randomized controlled trial & Remote physical activity program \\
\bottomrule
\end{longtable}

\noindent\textit{Note:} The active pool currently contains 34 primary studies. Final PRISMA study/report counts remain provisional pending complete search export and deduplication.

\section*{Table 2. Principal findings from higher-priority and quantitatively informative studies}

\begin{longtable}{p{0.15\textwidth} p{0.20\textwidth} p{0.31\textwidth} p{0.27\textwidth}}
\toprule
Study & Domain & Key result & Interpretation \\
\midrule
\endfirsthead
\toprule
Study & Domain & Key result & Interpretation \\
\midrule
\endhead
    Garbarino (2015) & Work stress & aOR 2.68 (95\% CI 1.08-6.70) & Higher work stress predicted incident MetS \\
    Kuwahara (2016) & Physical activity & HR 0.74 (95\% CI 0.62-0.89) for 7.5 to <16.5 MET-h/week; HR 0.81 (0.69-0.96) for $\geq$16.5 & Higher-intensity mixed leisure exercise associated with lower MetS incidence; commuting walking not associated \\
    Browne et al. (2017) & Physical activity & Active vs sedentary: MetS OR 0.52 (95\% CI 0.27-0.98); abdominal obesity OR 0.36 (0.16-0.82); elevated BP OR 0.47 (0.26-0.84); low HDL OR 0.54 (0.31-0.93) & Meeting PA recommendations was associated with lower MetS and component risk \\
    Yamaguchi (2018) & Work-related stress & aOR 3.27 (95\% CI 1.46-7.34) & Increasing job demands over time associated with higher incident MetS \\
    Haufe (2019) & Exercise & Between-group difference -0.26 (95\% CI -0.35 to -0.16); p<0.0001 & Telemonitoring-supported exercise reduced MetS severity \\
    Aboma (2020) & Lifestyle/diet/sedentary behaviour & Sedentary $\geq$8 h: APR 2.29 (95\% CI 1.88-2.80); $\geq$5 fruit/veg servings: APR 0.04 (0.01-0.27) & Sedentary time associated with higher MetS prevalence; high fruit/vegetable intake with lower prevalence \\
    van Zon et al. (2020) & Occupation/physical activity/diet & Men stationary plant/machine operators HR 1.94 (95\% CI 1.26-3.00); women food-preparation assistants HR 1.80 (1.01-3.22) & MetS risk differed by occupation and sex \\
    van Zon et al. (2021) & Diet/physical activity/health behaviours & Low-skilled white-collar OR 1.24 (95\% CI 1.12-1.37); low-skilled blue-collar OR 1.37 (1.18-1.59) & Older low-skilled workers had higher incident MetS; behaviours explained part of differences \\
    Eftekhari (2021) & Job stress & Office workers OR 1.51 (95\% CI 1.25-1.82); service personnel OR 1.74 (1.41-2.14) vs clinical staff & MetS prevalence varied by employment category; psychosocial well-being was poorer with MetS \\
    Zhang (2024) & Occupational stress & No overall association with MetS; adjusted models linked D/C ratio with SBP, social support with waist circumference, and over-commitment with fasting blood glucose & Occupational stress was not associated with MetS overall, but selected stress dimensions were associated with specific MetS components \\
    Sharifi (2025) & Diet quality & MetS prevalence 22.3\%; diet moderation associated with MetS OR 1.03 (95\% CI 1.008-1.05), abdominal obesity OR 1.02 (1.003-1.04), elevated TG OR 1.02 (1.00-1.03), elevated FBS OR 1.03 (1.005-1.05) & Diet moderation score was associated with MetS and several components; adequacy was not \\
    Ma (2025) & Dietary intervention & FBG beta -0.72 p=0.010; TC -1.49 p<0.001; LDL-C -0.65 p<0.001; WC -7.73 p<0.001; BMI -2.01 p<0.001; HDL-C +0.13 p<0.001 & Canteen-based healthy lunch plus advice improved several MetS components \\
    Ayeltigah et al. (military) (2026) & Work stress/diet/physical activity & Work stress and MetS: OR 6.472 (95\% CI 1.535-27.292) & Work stress associated with higher MetS odds; estimate imprecise \\
    Earnest (2015) & Worksite behavioral weight-loss intervention & Women 43\% to 30\%; men 52\% to 26\%; p<0.001 & Voluntary worksite weight-loss program reduced MetS prevalence and risk factors \\
    Ryu (2017) & Workplace health promotion: education/self-monitoring/diet/activity & Targeted intensive group: waist 89.96 to 86.93 cm and fasting glucose 93.44 to 84.56 mg/dL, p<0.001; MetS prevalence 41.5\% to 31.7\% & Targeted proactive workplace program improved MetS indicators; voluntary low-intensity components did not \\
    So (2025) & HIIT plus calorie restriction & Weight, blood pressure and triglycerides improved post-intervention and remained improved at 1 year (p<0.01) & Web-based HIIT + calorie restriction improved MetS-related outcomes \\
    Bogale (2025) & Lifestyle education: diet/physical activity/stress and other risk factors & Waist -5.33 cm; SBP -6.96 mmHg; DBP -4.21 mmHg; total cholesterol -34.12 mg/dL; LDL -20.68 mg/dL; all p<0.0001 & Workplace lifestyle education improved several MetS risk factors \\
    Kim et al. BEST (2015) & Physical activity/diet/counselling & Weight p=0.022; visceral fat p=0.033; waist circumference p=0.037 & Internet-based activity/diet/counselling improved selected cardiometabolic risks \\
    Choi (2017) & T'ai Chi + health education & SBP p=0.003; DBP p=0.037; triglycerides p=0.019 & T'ai Chi plus health education improved blood pressure and triglycerides \\
    Huang et al. (2017) & Occupational/leisure/commuting physical activity & High vs low leisure-time PA: OR 0.76 (95\% CI 0.62-0.95) & Higher leisure-time PA associated with lower MetS odds; commuting PA not associated \\
    Fang et al. (2019) & Physical activity program & Reduced number of MetS risk factors; lower triglycerides, total cholesterol and LDL & Worksite PA program improved metabolic and fitness outcomes \\
    Schilling et al. (2020) & Work stress/physical activity/cardiorespiratory fitness & Cardiorespiratory fitness predicting MetS risk: beta -0.38 & Higher fitness predicted lower follow-up MetS risk; no main effect for PA or work stress \\
    Park \& Hwang (2024) & Remote physical activity program & Intervention improved physical and physiological indicators versus control (P<0.05) & Remote workplace PA education improved MetS-related indicators \\
\bottomrule
\end{longtable}

\noindent\textit{Abbreviations:} APR, adjusted prevalence ratio; aOR, adjusted odds ratio; BMI, body mass index; CI, confidence interval; FBG, fasting blood glucose; HDL, high-density lipoprotein; HR, hazard ratio; LDL, low-density lipoprotein; MetS, metabolic syndrome; OR, odds ratio; PA, physical activity; SBP, systolic blood pressure; WC, waist circumference.

\section*{Table 3. Provisional methodological appraisal summary}

\begin{tabular}{p{0.24\textwidth} p{0.07\textwidth} p{0.18\textwidth} p{0.43\textwidth}}
\toprule
Design group & n & JBI tool & Main appraisal pattern \\
\midrule
    Randomized controlled trials & 6 & JBI RCT & Randomized comparisons strengthen causal inference, but concealment, masking, attrition and analysis by randomized assignment remain incompletely reported in some trials. \\
    Prospective/cohort studies & 6 & JBI Cohort & Generally stronger temporal ordering; attrition and missing-data handling are the main recurring uncertainties, and one police study did not require MetS absence at baseline. \\
    Analytical cross-sectional studies & 17 & JBI analytical cross-sectional & Useful for occupational associations but unable to establish temporality; reverse causation remains possible. \\
    Quasi-/nonrandomized interventions & 5 & JBI quasi-experimental & Useful intervention evidence but weaker causal confidence where randomization, repeated measurements, or concurrent controls are absent. \\
\bottomrule
\end{tabular}

\noindent\textit{Note:} All studies in the active pool have now undergone preliminary design-appropriate appraisal. No arbitrary summed quality score is used; appraisal findings qualify confidence in the narrative synthesis.
